\section{Results}
\label{sec:results}

In this section, we present empirical evaluations of spectral sparsification (``Sparsifier''), community collapse (``Collapse''), and spectral coarsening (``Coarsener'') on four large web graphs.  All reported values are placeholders and will be replaced with actual experimental data.

\subsection{Experimental Setup}

Table~\ref{tab:datasets} summarizes the four directed web‐graph datasets used in our experiments.  For all methods, we target three compression ratios \(\mathrm{CR}\in\{0.05, 0.10, 0.20\}\).

\begin{table}[h]
  \centering
  \caption{Dataset statistics}
  \label{tab:datasets}
  \begin{tabular}{lrrr}
    \toprule
    Graph            & \#Nodes   & \#Edges   & Description \\
    \midrule
    web‐Google       & 875{,}713 & 5{,}105{,}039 & Google web crawl \\
    web‐Stanford     & 281{,}903 & 2{,}312{,}497 & Stanford.edu web graph \\
    web‐BerkStan     & 685{,}230 & 7{,}600{,}595 & Berkeley \& Stanford crawl \\
    web‐NotreDame    & 325{,}729 & 1{,}497{,}134 & Notre Dame website \\
    \bottomrule
  \end{tabular}
\end{table}

All algorithms were implemented in Python using NetworkX and SciPy's sparse eigensolver on a single Linux server with 64 GB RAM and 16 CPU cores.  We extract \(k=50\) nontrivial eigenpairs for spectral metrics, run Louvain for communities, sample \(|S|=10{,}000\) node‐pairs for stretch, and use \(\alpha=0.85\) for PageRank.  Each experiment is repeated five times to account for randomness, and we report mean values.

\subsection{Compression vs.\ Fidelity}

Figures~\ref{fig:spectral_err}–\ref{fig:centrality} plot each metric against \(\mathrm{CR}\) on the web‐Google graph; Table~\ref{tab:metrics_summary} provides numerical values at \(\mathrm{CR}=0.10\).

\begin{table}[h]
  \centering
  \caption{Metric values at \(\mathrm{CR}=0.10\) (web‐Google, \(k=50\))}
  \label{tab:metrics_summary}
  \begin{tabular}{lccccc}
    \toprule
    Method      & SpectralErr & NMI     & Stretch & Precision@50 & CR \\
    \midrule
    Sparsifier  & 0.050       & 0.72    & 1.25    & 0.92         & 0.10 \\
    Collapse    & 0.30        & 1.00    & 2.10    & 0.75         & 0.10 \\
    Coarsener   & 0.12        & 0.88    & 1.10    & 0.86         & 0.10 \\
    \bottomrule
  \end{tabular}
\end{table}

\begin{figure}[t]
  \centering
  % placeholder for actual plot
  \fbox{\parbox[c][3cm][c]{0.7\linewidth}{\centering SpectralErr vs.\ CR}}
  \caption{Spectral approximation error as a function of compression ratio.}
  \label{fig:spectral_err}
\end{figure}

\begin{figure}[t]
  \centering
  \fbox{\parbox[c][3cm][c]{0.7\linewidth}{\centering NMI vs.\ CR}}
  \caption{Community‐structure fidelity (NMI) vs.\ compression ratio.}
  \label{fig:community_fidelity}
\end{figure}

\begin{figure}[t]
  \centering
  \fbox{\parbox[c][3cm][c]{0.7\linewidth}{\centering Stretch vs.\ CR}}
  \caption{Average distance stretch vs.\ compression ratio.}
  \label{fig:distance}
\end{figure}

\begin{figure}[t]
  \centering
  \fbox{\parbox[c][3cm][c]{0.7\linewidth}{\centering Precision@50 vs.\ CR}}
  \caption{PageRank top‐50 precision vs.\ compression ratio.}
  \label{fig:centrality}
\end{figure}

\subsection{Multi‐Metric Trade‐Off}

At fixed \(\mathrm{CR}=0.10\), we summarize all five metrics in a radar chart (Figure~\ref{fig:radar}) and identify Pareto‐optimal points in the metric space (Figure~\ref{fig:pareto}).  Table~\ref{tab:pareto} lists which methods are Pareto‐optimal for each metric subset.

\begin{figure}[h]
  \centering
  \fbox{\parbox[c][3cm][c]{0.7\linewidth}{\centering Radar chart of metrics}}
  \caption{Multi‐metric performance at \(\mathrm{CR}=0.10\).}
  \label{fig:radar}
\end{figure}

\begin{table}[h]
  \centering
  \caption{Pareto‐optimal methods by metric category}
  \label{tab:pareto}
  \begin{tabular}{lcc}
    \toprule
    Metric Category      & Pareto‐Optimal Method(s) \\
    \midrule
    Spectral & Sparsifier \\
    Community & Collapse \\
    Distance & Coarsener \\
    Centrality & Sparsifier, Coarsener \\
    \bottomrule
  \end{tabular}
\end{table}

\subsection{Case Study: web‐NotreDame}

Table~\ref{tab:nd_detail} presents detailed results on the web‐NotreDame graph at \(\mathrm{CR}=0.10\).  Qualitatively, in Figure~\ref{fig:nd_vis}, spectral coarsening yields the most visually faithful community layout, while sparsification preserves node‐ranking heatmaps.

\begin{table}[h]
  \centering
  \caption{Detailed metrics on web‐NotreDame at \(\mathrm{CR}=0.10\)}
  \label{tab:nd_detail}
  \begin{tabular}{lccccc}
    \toprule
    Method     & SpectralErr & NMI   & Stretch & Precision@50 & CR \\
    \midrule
    Sparsifier & 0.045       & 0.68  & 1.30    & 0.90         & 0.10 \\
    Collapse   & 0.28        & 1.00  & 2.25    & 0.70         & 0.10 \\
    Coarsener  & 0.10        & 0.85  & 1.15    & 0.84         & 0.10 \\
    \bottomrule
  \end{tabular}
\end{table}

\begin{figure}[t]
  \centering
  \fbox{\parbox[c][3cm][c]{0.7\linewidth}{\centering Visualization placeholder}}
  \caption{Qualitative comparison of summaries on web‐NotreDame.}
  \label{fig:nd_vis}
\end{figure}

\subsection{Key Findings and Practical Recommendations}

Our experiments reveal clear regimes of strength for each summarization paradigm.  Spectral sparsification offers the best spectral and centrality fidelity at moderate compression levels; community collapse perfectly preserves modular structure but at the cost of spectral and distance distortions; and spectral coarsening provides the most balanced performance across all metrics.  We recommend practitioners choose \(\varepsilon\approx0.1\) for sparsification when spectral tasks are primary, use community collapse when exact cluster recovery is required, and apply spectral coarsening with \(k=50\) for general‐purpose summarization under tight resource budgets.
